\week{W4}{Parallel manipulators}{4}
\begin{multicols}{3}
\hd{Serial vs parallel \pg{9--19, 74}}
Parallel: “end effector connected to the base by multiple parallel kinematic chains” \pg{12}. More joints than DoF, so some joints are passive. Examples: Gough tyre tester, Stewart flight simulator (1965), ABB FlexPicker delta, surgical hexapod.
\begin{tabularx}{\linewidth}{@{}Xll@{}}\toprule
 & Serial & Parallel\\\midrule
Chains & open & closed\\
Joints & all active & active + passive\\
FK & simple & \textbf{complex}\\
IK & complex & \textbf{simple}\\
Joint errors & accumulate & do not\\
Workspace & large & small\\
Payload/weight & low & very high\\
Moving inertia & high & low\\\bottomrule
\end{tabularx}
\use{scenario “which architecture?”: fast pick-and-place $\to$ delta; heavy simulator $\to$ Stewart; precise surgical tool $\to$ hexapod; justify by stiffness, low moving mass, non-accumulating errors, and note the small workspace.}

\hd{Workspace of parallel robots \pg{20--23}}
Four kinds \pg{21}: \textbf{reachable} (some orientation), \textbf{dexterous} (all), \textbf{constant-orientation} (points reachable with one given orientation), \textbf{orientation workspace} (all orientations at one given point). Compute by \emph{discretisation} (step each actuator, take the boundary; slow, accuracy set by step) or \emph{geometry} (intersect the regions each actuator allows) \pg{22--23}.\\
Two-leg planar, prismatic legs: each leg confines the point to an annulus $r_{min}\le r\le r_{max}$ about its base pivot; \textbf{workspace = overlap of the annuli}. Sketch: draw each annulus, shade the intersection.

\hd{Mobility: G--K \pg{24--34}}
\begin{fbx}[Plain Grübler--Kutzbach \pg{25}]
\[M=D\,(N-J-1)+\sum_{i=1}^{J}F_i\]
$D=3$ planar, $6$ spatial; $N$ includes ground; $F_R=F_P=1$, $F_U=2$, $F_S=3$. Derivation: $D(N{-}1)$ free DoF, each joint removes $D-F_i$.
\use{first estimate for any mechanism.}
\end{fbx}
\begin{recipe}[Mobility of a parallel robot (W13 Q6, Q7)]
\label{w4:gk}
\begin{enumerate}
\item Planar or spatial? Motion out of the page means $D=6$ (“still 3D despite this depiction”, W13 p.14).
\item Number every rigid body, \textbf{ground = 1}. A prismatic pair is two bodies.
\item List every joint and its $F_i$. Three bodies at one pin $=$ two joints.
\item Substitute. If $M$ disagrees with the obvious motion, look for idle spins ($R_M$) or repeated constraints ($R_C$).
\end{enumerate}
\end{recipe}

\columnbreak
\begin{tabularx}{\linewidth}{@{}Xcccrr@{}}\toprule
Mechanism & $D$ & $N$ & $J$ & $\sum F$ & $M$\\\midrule
3-RPR \pg{26} & 3 & 8 & 9 & 9 & 3\\
Delta R--U--U legs \pg{30} & 6 & 8 & 9 & 15 & 3\\
Delta, parallelogram \pg{33} & 6 & 11 & 15 & 39 & 9 (!)\\
Four-bar \added & 3 & 4 & 4 & 4 & 1\\
Gough--Stewart 6-UPS \added & 6 & 14 & 18 & 36 & 6\\\bottomrule
\end{tabularx}
Delta R--U--U: $\sum F=(1+2+2)\cdot3=15$, $6(8-9-1)+15=3$. Common delta: $\sum F=(1+4\cdot3)\cdot3=39$, $6(11-15-1)+39=9$. Stewart: $6(14-18-1)+36=6$. 3-RPR: $3(8-9-1)+9=3$. \chk\\
The simplest delta is not built: stiff, play-free universal joints are hard to make and play lowers the usable frequency \pg{31}.
\begin{slideerr}[slide 29]
A handwritten note reads $D(N-J+1)$. The printed formula and every worked count use $D(N-J-1)$. With $+1$ the simple delta would give $6(8-9+1)+15=15$. \chk
\end{slideerr}
\begin{trap}
Slider-crank (ground, crank, rod, slider; 3 R + 1 P) with $D=6$: $6(4-4-1)+4=-2$, “a structure”. It is planar: $3(4-4-1)+4=1$. Spatial example: $N=6$, $J=7$ (5 R, 1 U, 1 S): $\sum F=10$, $M=6(6-7-1)+10=-2$: report it, then look for $R_C$ before saying “structure”. \chk
\end{trap}

\hd{Extended G--K \pg{35--43}}
\begin{fbx}[Equation \pg{35}]
\[M=D(N-J-1)+\sum_{i=1}^{J}F_i+R_C-R_M\]
\[R_C=\sum_{i=1}^{J}F_i-\sum_{k=1}^{L}u_k\]
$R_M$ = \textbf{local mobility}: internal motion that does not move the output (a rod with a ball joint at each end spins about its own axis). Plain G--K counts it, so \textbf{subtract}. $R_C$ = \textbf{over-constraints}: joints repeating a constraint, so freedoms are removed twice: \textbf{add back}. $L$ = number of loops; $u_k$ = distinct platform DoF loop $k$ allows (each distinct DoF counted once across loops).
\use{plain $M$ implausible (delta 9; 4RRR $-6$). Say \emph{which} of $R_C$, $R_M$ applies and why.}
\end{fbx}
\begin{worked}[common delta \pg{33--36}]
Plain 9, but the platform only translates (3 DoF). Each of the 6 rods spins about its own axis: $R_M=6$. No repeated constraints: $R_C=0$. $M=9+0-6=\mathbf 3$. \chk
\end{worked}
\begin{worked}[4RRR \pg{37--43}]
Two square plates, four RR legs: $N=10$, $J=12$, $\sum F=12$. Plain: $6(10-12-1)+12=-6$ (impossible). Loop 1 allows $\{T_y,T_z,R_x\}$, loop 2 $\{T_x,T_z,R_y\}$; union $\{T_x,T_y,T_z,R_x,R_y\}$, $\sum u_k=5$ ($T_z$ once). $R_C=12-5=7$, $R_M=0$: $M=-6+7-0=\mathbf 1$ (platform only moves in $T_z$). \chk
\end{worked}
“Explain” marks: plain G--K assumes every constraint is independent ($R_C$ fixes that) and every DoF moves the output ($R_M$ fixes that).

\columnbreak
\hd{Active, passive joints; loop closure \pg{45--58}}
Configuration $\mathbf q=(\boldsymbol\theta,\boldsymbol\varphi)$: $n$ active $\boldsymbol\theta$, $m$ passive $\boldsymbol\varphi$ \pg{45}. Gough--Stewart: 6 active, 30 passive variables \pg{58}.
\begin{recipe}[Loop-closure equations \pg{49}]
\label{w4:loop}
\begin{enumerate}
\item Break the loop into serial chains (choose the break giving the fewest equations).
\item Write the pose of the break point from each side.
\item Equate position and orientation. Count unknowns: active $\theta$, passive $\varphi$.
\end{enumerate}
\end{recipe}
\begin{fbx}[Four-bar, break link 2 at $K$ \pg{54}]
$a,b$ = distances from $K$ to link 2’s ends ($L_2=a+b$); {D} at $(L_0,0)$:
\[\begin{aligned}L_1c_{\theta_1}+a\cos(\theta_1{+}\varphi_2)&=L_0+L_3c_{\varphi_1}+b\cos(\varphi_1{+}\varphi_3)\\ L_1s_{\theta_1}+a\sin(\theta_1{+}\varphi_2)&=L_3s_{\varphi_1}+b\sin(\varphi_1{+}\varphi_3)\\ \theta_1+\varphi_2&=\varphi_1+\varphi_3+\pi\end{aligned}\]
Three equations: 1 active ($\theta_1$), 3 passive ($\varphi_1,\varphi_2,\varphi_3$).
\end{fbx}
\begin{fbx}[Freudenstein, break at both ends \pg{57}]
\[\eta_1=(L_1c_{\theta_1}-L_0-L_3c_{\varphi_1})^2+(L_1s_{\theta_1}-L_3s_{\varphi_1})^2-L_2^2=0\]
\use{only link 2’s length must be preserved: 1 equation in $\theta_1,\varphi_1$. Given $\theta_1$, two $\varphi_1$ (two assembly modes).}
\end{fbx}

\hd{Forward kinematics (hard) \pg{59--67}}
Given active joints, solve loop closure for the passive ones, then the pose. No general method \pg{60}; solutions can be multiple or imaginary, so numerics are common \pg{67}.\\
\textbf{Four-bar by geometry} \added: $A=L_1(\cos\theta_1,\sin\theta_1)$ and $D=(L_0,0)$ fixed; $B$ is on the circle $L_2$ about $A$ and $L_3$ about $D$. Two intersections = two \textbf{assembly modes}; one = limit; none = cannot assemble at that $\theta_1$.
\begin{fbx}[Slide route \pg{62--66}: from $K=(x,y)$ find $\delta=\theta_1+\varphi_2$]
$A_1=x-L_0,\ B_1=y,\ C_1=\dfrac{(x-L_0)^2+y^2+b^2-L_3^2}{2b}$\\
$A_2=x,\ B_2=y,\ C_2=\dfrac{L_1^2-a^2-x^2-y^2}{2a}$\\
$A_i\cos\delta+B_i\sin\delta+C_i=0$. Cramer:
\[\cos\delta=\frac{B_1C_2-B_2C_1}{A_1B_2-A_2B_1},\ \ \sin\delta=\frac{A_2C_1-A_1C_2}{A_1B_2-A_2B_1}\]
$\delta=\atantwo(\sin\delta,\cos\delta)$. Consistency (slide 65): $(A_1B_2-A_2B_1)^2=(A_1C_2-A_2C_1)^2+(B_1C_2-B_2C_1)^2$.
\end{fbx}
\begin{slideerr}[slide 65]
Printed:
\[\delta=\tan^{-1}\frac{A_1C_2-A_2C_1}{(B_1C_2-B_2C_1)+(A_1B_2-A_2B_1)}\]
Tangent-half-angle gives
\[\delta=2\tan^{-1}\frac{A_2C_1-A_1C_2}{(B_1C_2-B_2C_1)+(A_1B_2-A_2B_1)}\]
opposite sign and a factor 2. Test: slide form $0.300^\circ$, true $-0.600^\circ$. Use the atan2 of the Cramer sine and cosine. \chk
\end{slideerr}

\columnbreak
\begin{worked}[four-bar FK, same as Sample Q5]
$L_0=1.0$, $L_1=0.5$, $L_3=0.6$, $a=0.2$, $b=0.25$ ($L_2=0.45$), $\theta_1=70^\circ$:
\begin{tabularx}{\linewidth}{@{}Xccc@{}}\toprule
Mode & $K$ (m) & $\varphi$ & $\varphi_1$\\\midrule
$B$ above $AD$ & (0.3710, 0.4678) & $-0.600^\circ$ & $129.175^\circ$\\
$B$ below $AD$ & (0.2756, 0.2993) & $-58.486^\circ$ & $171.739^\circ$\\\bottomrule
\end{tabularx}
Below mode by hand: $A=(0.1710,0.4698)$, $|AD|=0.9529$. $x=(L_2^2-L_3^2+|AD|^2)/(2|AD|)=(0.2025-0.36+0.9080)/1.9058=0.3938$, $h=\sqrt{0.45^2-0.3938^2}=0.2178$. $B=(0.4062,0.0862)$, $\varphi=-58.486^\circ$, $K=A+0.2(\cos\varphi,\sin\varphi)=(0.2756,0.2993)$. \chk\\
Non-Grashof ($0.45+1.0>0.5+0.6$): no link turns fully; assembles only for $|\theta_1|\le81.52^\circ$ ($|AD|\le L_2+L_3$).
\end{worked}

\hd{Inverse kinematics (easy) \pg{68--73}}
Given the output pose, each leg is a serial chain with known end pose: solve each leg alone, no loop closure \pg{69}. A redundant leg gives a non-unique answer.
\begin{recipe}[Four-bar IK, from $K=(x,y)$ and $\varphi$ (W13 Q5) \pg{70--72}]
\label{w4:fourbar}
\begin{enumerate}
\item Link 2 ends: $A=K-a(\cos\varphi,\sin\varphi)$, $B=K+b(\cos\varphi,\sin\varphi)$.
\item $\theta_1=\atantwo(y-a\sin\varphi,\ x-a\cos\varphi)$.
\item $\varphi_2=\varphi-\theta_1$.
\item $\varphi_1=\atantwo(y+b\sin\varphi,\ x+b\cos\varphi-L_0)$.
\item $\varphi_3=\varphi+\pi-\varphi_1$ (mod $2\pi$).
\item Check: with the table above $K=(0.3710,0.4678)$, $\varphi=-0.600^\circ$ returns $\theta_1=70^\circ$, $\varphi_1=129.175^\circ$.
\end{enumerate}
\end{recipe}
\begin{fbx}[Gough--Stewart leg (homework \pg{73})]
Leg $i$: base point $\mathbf b_i$ in {B}, platform point $\mathbf p_i$ in {P}, ${}^B\T_P=[\Rm\,|\,\mathbf t]$:
\[\mathbf l_i=\mathbf t+\Rm\,\mathbf p_i-\mathbf b_i,\quad d_i=\lVert\mathbf l_i\rVert,\quad\hat{\mathbf n}_i=\mathbf l_i/d_i\]
$d_i$ = active prismatic length. Base S takes the azimuth and elevation of $\hat{\mathbf n}_i$; platform U takes the angle between $\hat{\mathbf n}_i$ and the platform axes $\Rm$. One answer per leg.
\end{fbx}
\begin{worked}[Stewart leg lengths]
Base points on a 1.0\,m circle at $-10,10,110,130,230,250^\circ$; platform points on a 0.5\,m circle at $-50,50,70,170,190,290^\circ$. Home (level, 1.0\,m up), leg 1:
\[\mathbf l_1=(-0.6634,-0.2094,1.0),\quad d_1=\sqrt{0.4401+0.0438+1}=1.2182\text{ m}\]
All six legs $=1.2182$\,m (leg 1 azimuth $-162.48^\circ$, elevation $55.17^\circ$).\\
Pose $t(0.05,-0.02,1.0)\,R_z(5^\circ)R_y(3^\circ)$: legs (m)
\[1.1597,\ 1.1974,\ 1.2170,\ 1.2812,\ 1.2152,\ 1.2504\ \chk\]
\end{worked}
Delta IK (demo, base radius 0.20, platform 0.05, arm 0.20, rods 0.45\,m): each arm alone, like the four-bar; at $\mathbf p=(0.03,-0.02,-0.38)$ the motor angles are $16.339^\circ,\,30.871^\circ,\,22.996^\circ$ and every rod stays $0.45$\,m (that is the loop closure). Parallelograms keep the platform level: 3 translations, no rotation. \chk
\begin{trap}
FK/IK swapped: FK is the hard direction for parallel robots (coupled nonlinear loops), IK the easy one. For serial arms it is the opposite.
\end{trap}
\textbf{Practical:} none this week.
\end{multicols}
